\begin{afterword}
\summaryhead

The post opens with a passage from Robyn Dawes. In experiments where 70\% of the cards are blue and the order is random, the best strategy is to predict blue every time, which is right 70\% of the time. Subjects instead predict blue about as often as it appears, which is right only 58\% of the time; even when paid a nickel for each correct guess over a thousand trials, they predicted blue 76\% of the time.

The author argues that subjects could have tested their hunches about patterns privately while still betting blue. The author suspects that the all-blue strategy simply does not occur to people, who expect the best strategy to look like the random sequence. The post calls it a counterintuitive idea that one should act and think lawfully under uncertainty: randomizing one's actions only takes one further from the target. Most people fight chaos with chaos, which is why there are few rationalists, and every departure from lawful thinking is an expected loss.

\responsehead

The post's reasoning about the card task is sound. When 70\% of the cards are blue and the order is random, betting blue every time is best, and the post's arithmetic for the alternatives is right. Its suspicion that the all-blue strategy does not occur to people was offered as a suspicion, and later experiments support it. The problems are in the general claims the post draws from the task.

The general claims go well beyond the task. From one laboratory result the post concludes that ``most who perceive a chaotic world will try to fight chaos with chaos,'' and that this is why there are few rationalists. But the behaviour it generalizes from changes with practice and pay. In a 1000-trial study from 1961, subjects ended up choosing the likelier option 83\% of the time, and in experiments published in 2002, about 70\% of subjects who were well paid, regularly informed and long trained came to choose it on every trial for at least fifty trials in a row. That the task shows ``the most important idea in all of rationality'' is asserted, not argued.

The rule that randomizing ``takes you further from the target'' is stated without its exception. Against an opponent who can predict your method, choosing at random is the best one can do, and the author granted this two days later, in a post this book does not include. This post turns to opponents but considers only an ``irrational'' one. A reader of the book meets the rule without the case where it fails.

In short: a sound point about one betting task, extended to claims about most people and about all uncertainty, with the exception to its rule left to a later post that the book omits.
\end{afterword}
